Let $F$ be the forget set and $R$ the retain set.
An unlearning method starts from a model $\theta_{\mathrm{full}}$ trained on both sets and produces an unlearned model $\theta_0$.
The goal is not only that $\theta_0$ fails to recover answers from $F$ immediately, but that facts in $F$ remain suppressed after subsequent relearning updates.

For each forget item $i \in F$, the evaluator uses prompt variants $p_{i,j}$, a target answer $a_i$, optional aliases, and a neutral answer $n_i$.
The primary score is a target log-likelihood margin,
\[
  m_i(\theta) = \max_j \left[\log p_{\theta}(a_i \mid p_{i,j}) - \log p_{\theta}(n_i \mid p_{i,j})\right].
\]
Thresholds $\tau_i$ are calibrated from a retain-only oracle using fixed quantiles.
An item is resurrected under the margin metric when
\[
  m_i(\theta) > \tau_i.
\]

After unlearning, the evaluator applies relearning updates on a pool $B$ for horizons
\[
  t \in \{0,1,2,4,8,16,32,64,128,256,512\}.
\]
The model after $t$ relearning steps is $\theta_t$.
The basic resurrection fraction is
\[
  r_t = \frac{1}{|F|}\sum_{i \in F} \mathbf{1}[m_i(\theta_t) > \tau_i],
\]
and survival is $S_t = 1-r_t$.

This setup separates three different requirements:

\begin{enumerate}[leftmargin=*]
  \item immediate forgetting: low $r_0$;
  \item retained utility: low retain negative log-likelihood and low refusal;
  \item durability: low resurrection growth under relearning.
\end{enumerate}
